\documentclass{article}
\usepackage{amsmath,amssymb}
\usepackage{xurl}
\usepackage[hidelinks]{hyperref}
\setlength{\emergencystretch}{2em}
% Deferred to a later review branch: % Shared commands for the notes in docs/paper-gaps/.
% Every note loads this file via % Shared commands for the notes in docs/paper-gaps/.
% Every note loads this file via % Shared commands for the notes in docs/paper-gaps/.
% Every note loads this file via \input{command}. Project-wide notation
% belongs here, never in per-note redefinitions.

% --- Project configuration ------------------------------------------------
\newcommand{\projectIssueBase}{https://github.com/LionSR/QICLean/issues/}
\newcommand{\projectPrBase}{https://github.com/LionSR/QICLean/pull/}
\newcommand{\projectDocBase}{https://sirui-lu.com/QICLean/docs/}
% ---------------------------------------------------------------------------

\DeclareUnicodeCharacter{2080}{\ifmmode{}_{0}\else\textsubscript{0}\fi}
\DeclareUnicodeCharacter{2081}{\ifmmode{}_{1}\else\textsubscript{1}\fi}
\DeclareUnicodeCharacter{2082}{\ifmmode{}_{2}\else\textsubscript{2}\fi}
\DeclareUnicodeCharacter{2099}{\ifmmode{}_{n}\else\textsubscript{n}\fi}

\newcommand{\Lean}{\textsc{Lean}}

% A formal declaration name, upright sans, breakable at underscores.
\ExplSyntaxOn
\NewDocumentCommand{\leanid}{m}{
  \ifmmode
    \textsf{\detokenize{#1}}
  \else
    \group_begin:
    \urlstyle{sf}
    \tl_set:Nn \l_tmpa_tl {#1}
    \tl_replace_all:Nnn \l_tmpa_tl {\_} {_}
    \exp_args:NV \path \l_tmpa_tl
    \group_end:
  \fi
}
\ExplSyntaxOff

% Traceability links; use in footnotes, not in the running prose.
\newcommand{\ghissue}[1]{\href{\projectIssueBase#1}{GitHub issue \##1}}
\newcommand{\ghpr}[1]{\href{\projectPrBase#1}{PR \##1}}
\newcommand{\doclink}[1]{\href{\projectDocBase#1.html}{\path{#1}}}

% Domain notation, using \providecommand so a note-local refinement stays
% possible.
\providecommand{\tr}{\operatorname{tr}}

% Dirac notation and the identity operator, as in blueprint/src/macros.
\usepackage{dsfont}
\providecommand{\ket}[1]{|#1\rangle}
\providecommand{\bra}[1]{\langle #1|}
\providecommand{\braket}[2]{\langle #1 | #2 \rangle}
\providecommand{\ketbra}[2]{|#1\rangle\!\langle #2|}
\providecommand{\Id}{\mathds{1}}
\providecommand{\one}{\mathds{1}}
\providecommand{\C}{\mathbb{C}}
\providecommand{\MN}[1]{M_{#1}(\C)}
\providecommand{\MD}{M_D(\C)}

% Verdict marker: \gapnote{<kind>}{<status>}. Kind is the policy
% classification (clarification, local-correction, scope-restriction,
% unfaithful, false-source, open-gap); status is open, wip (elimination
% underway), resolved, or historical. Machine-read by the site index;
% prints a verdict line.
\newcommand{\gapnote}[2]{\par\noindent\textbf{Verdict:} #1 (#2).\par}
. Project-wide notation
% belongs here, never in per-note redefinitions.

% --- Project configuration ------------------------------------------------
\newcommand{\projectIssueBase}{https://github.com/LionSR/QICLean/issues/}
\newcommand{\projectPrBase}{https://github.com/LionSR/QICLean/pull/}
\newcommand{\projectDocBase}{https://sirui-lu.com/QICLean/docs/}
% ---------------------------------------------------------------------------

\DeclareUnicodeCharacter{2080}{\ifmmode{}_{0}\else\textsubscript{0}\fi}
\DeclareUnicodeCharacter{2081}{\ifmmode{}_{1}\else\textsubscript{1}\fi}
\DeclareUnicodeCharacter{2082}{\ifmmode{}_{2}\else\textsubscript{2}\fi}
\DeclareUnicodeCharacter{2099}{\ifmmode{}_{n}\else\textsubscript{n}\fi}

\newcommand{\Lean}{\textsc{Lean}}

% A formal declaration name, upright sans, breakable at underscores.
\ExplSyntaxOn
\NewDocumentCommand{\leanid}{m}{
  \ifmmode
    \textsf{\detokenize{#1}}
  \else
    \group_begin:
    \urlstyle{sf}
    \tl_set:Nn \l_tmpa_tl {#1}
    \tl_replace_all:Nnn \l_tmpa_tl {\_} {_}
    \exp_args:NV \path \l_tmpa_tl
    \group_end:
  \fi
}
\ExplSyntaxOff

% Traceability links; use in footnotes, not in the running prose.
\newcommand{\ghissue}[1]{\href{\projectIssueBase#1}{GitHub issue \##1}}
\newcommand{\ghpr}[1]{\href{\projectPrBase#1}{PR \##1}}
\newcommand{\doclink}[1]{\href{\projectDocBase#1.html}{\path{#1}}}

% Domain notation, using \providecommand so a note-local refinement stays
% possible.
\providecommand{\tr}{\operatorname{tr}}

% Dirac notation and the identity operator, as in blueprint/src/macros.
\usepackage{dsfont}
\providecommand{\ket}[1]{|#1\rangle}
\providecommand{\bra}[1]{\langle #1|}
\providecommand{\braket}[2]{\langle #1 | #2 \rangle}
\providecommand{\ketbra}[2]{|#1\rangle\!\langle #2|}
\providecommand{\Id}{\mathds{1}}
\providecommand{\one}{\mathds{1}}
\providecommand{\C}{\mathbb{C}}
\providecommand{\MN}[1]{M_{#1}(\C)}
\providecommand{\MD}{M_D(\C)}

% Verdict marker: \gapnote{<kind>}{<status>}. Kind is the policy
% classification (clarification, local-correction, scope-restriction,
% unfaithful, false-source, open-gap); status is open, wip (elimination
% underway), resolved, or historical. Machine-read by the site index;
% prints a verdict line.
\newcommand{\gapnote}[2]{\par\noindent\textbf{Verdict:} #1 (#2).\par}
. Project-wide notation
% belongs here, never in per-note redefinitions.

% --- Project configuration ------------------------------------------------
\newcommand{\projectIssueBase}{https://github.com/LionSR/QICLean/issues/}
\newcommand{\projectPrBase}{https://github.com/LionSR/QICLean/pull/}
\newcommand{\projectDocBase}{https://sirui-lu.com/QICLean/docs/}
% ---------------------------------------------------------------------------

\DeclareUnicodeCharacter{2080}{\ifmmode{}_{0}\else\textsubscript{0}\fi}
\DeclareUnicodeCharacter{2081}{\ifmmode{}_{1}\else\textsubscript{1}\fi}
\DeclareUnicodeCharacter{2082}{\ifmmode{}_{2}\else\textsubscript{2}\fi}
\DeclareUnicodeCharacter{2099}{\ifmmode{}_{n}\else\textsubscript{n}\fi}

\newcommand{\Lean}{\textsc{Lean}}

% A formal declaration name, upright sans, breakable at underscores.
\ExplSyntaxOn
\NewDocumentCommand{\leanid}{m}{
  \ifmmode
    \textsf{\detokenize{#1}}
  \else
    \group_begin:
    \urlstyle{sf}
    \tl_set:Nn \l_tmpa_tl {#1}
    \tl_replace_all:Nnn \l_tmpa_tl {\_} {_}
    \exp_args:NV \path \l_tmpa_tl
    \group_end:
  \fi
}
\ExplSyntaxOff

% Traceability links; use in footnotes, not in the running prose.
\newcommand{\ghissue}[1]{\href{\projectIssueBase#1}{GitHub issue \##1}}
\newcommand{\ghpr}[1]{\href{\projectPrBase#1}{PR \##1}}
\newcommand{\doclink}[1]{\href{\projectDocBase#1.html}{\path{#1}}}

% Domain notation, using \providecommand so a note-local refinement stays
% possible.
\providecommand{\tr}{\operatorname{tr}}

% Dirac notation and the identity operator, as in blueprint/src/macros.
\usepackage{dsfont}
\providecommand{\ket}[1]{|#1\rangle}
\providecommand{\bra}[1]{\langle #1|}
\providecommand{\braket}[2]{\langle #1 | #2 \rangle}
\providecommand{\ketbra}[2]{|#1\rangle\!\langle #2|}
\providecommand{\Id}{\mathds{1}}
\providecommand{\one}{\mathds{1}}
\providecommand{\C}{\mathbb{C}}
\providecommand{\MN}[1]{M_{#1}(\C)}
\providecommand{\MD}{M_D(\C)}

% Verdict marker: \gapnote{<kind>}{<status>}. Kind is the policy
% classification (clarification, local-correction, scope-restriction,
% unfaithful, false-source, open-gap); status is open, wip (elimination
% underway), resolved, or historical. Machine-read by the site index;
% prints a verdict line.
\newcommand{\gapnote}[2]{\par\noindent\textbf{Verdict:} #1 (#2).\par}

\title{Simultaneous parameters for a soft improving family}
\author{}
\date{2026-11-03}
\begin{document}
\maketitle
\gapnote{scope-restriction}{open}

The efficient-family argument of Arad--Landau--Vazirani--Vidick uses a
fourth-order soft truncation and specifies the window and filter scales
in terms of the gap; see \cite[lines 1373--1420]{Arad2016RigorousRG}.
The result below is a different sufficient scalar calibration. It
makes no assertion about the printed gap dependence or running time.

\section{A simultaneous choice}

Fix a local dimension $q\ge2$, a positive gap $g$, and natural exponents
$a,b,c$. For an integer $u\ge2$, put
\begin{align}
    \ell&=u^8,\quad m=u^{10},\quad t=3m,\quad k=u^{16},\quad T=u^9,\\
    D&=(mk+1)(qk)^{36\ell+4\lfloor k/\ell\rfloor},\\
    M&=3t(1+\log m)+4\ell+34,\\
    \tau&=D^{-a}M^{-b},\qquad
    \xi=\min\left\{1,\frac{\tau g^2}{20000(M+g+8)}\right\}.
\end{align}
The count and precision are evaluated at the same $u$. In particular,
one does not first choose a residual tolerance and then allow the
resulting family count to change that tolerance.

There is a sufficiently large natural $u$, depending only on
$q,g,a,b,c$, for which $t,T\ge156+2g$, $T\ge g/40+24$, and
\begin{align}
    3\left[32e^{-(T-g/40-24)/8}
       +\frac tm\left(\frac Tt\right)^m\right]
       &\le\frac{\tau g}{10^6},\\
    3\left[32e^{-(T-156-2g)/8}
       +\frac tm\left(\frac Tt\right)^m\right]
       &\le\frac{\tau g}{10^6},\\
    \left[2e^{-2k\sqrt{(g/2)/(M+\xi+4+g)}}\right]^2
       &\le D^{-12}M^{-c}.
\end{align}
Both $\tau$ and $\xi$ lie in $(0,1]$. This is a proved numerical
statement.\footnote{Formal theorem:
\leanid{Real.exists_soft_family_calibration} in
\doclink{QICLean/Analysis/SoftFamilyCalibration}.}

Indeed, $k/\ell=u^8$, and the exact expressions satisfy
\begin{align}
    \log D&\le u^8(\log2+40\log q+666\log u),\\
    M+R&\le u^{10}(47+R+90\log u)\qquad(R\ge0).
\end{align}
Since $\log u=o(u^p)$ for every $p>0$, eventually
$M+5+g\le u^{12}$ and
\begin{align}
    (a+12)\log D+(b+c)\log M
      +200+g+\left|\log\frac{g}{2000000}\right|
      \le\frac{u^9}{64}.
\end{align}
These inequalities control the two exponential residuals and the
Taylor remainder $9(3u)^{-m}$. They also control the squared filter
error, because, once $gu^2\ge1$,
\begin{align}
    u^9\le4u^{16}\sqrt{\frac{g/2}{M+\xi+4+g}}.
\end{align}

\section{The algorithmic distinction}

This calibration permits a precision depending on the family count,
with all choices made before the chain length or target dimension.
It does not complete a particular randomized construction. One possible
use is a procedure that repeatedly samples half the current dimension,
applies an improving family, transfers viability to the original
physical target, and then trims against that target. This differs from
the sample-once iteration of Figure~\texttt{fig:merge-prime}
in \cite[lines 1664--1673]{Arad2016RigorousRG}; its actual sampling law,
success probability, dimension decrease and trimming estimates must be
proved separately. No conclusion about that iteration follows solely
from the numerical existence theorem.

For a supplied common improving family and a reference satisfying the
proved fixed-input sampling and collective-trimming estimates, the actual
stopped half-resampling iteration now has a finite-product original-target
probability bound. One coefficient $K$, chosen before the length $N$,
target dimension bound $r$ and ambient dimension, gives stopping threshold
$K(r+1)(N+2)^2$. After $R$ prescribed rounds, the true event that every
actual state is $0.04$-viable, and each active next state is
$3/200$-viable, has failure probability at most $R\eta/(N+2)^a$.
The initial matrix need not be an orthogonal projection. The proof uses
actual guarded failures, the preceding-seed dependence of the recursion,
and fixed-input Haar estimates; no measurable selection of auxiliary
success events is assumed.\footnote{Formal theorem:
\leanid{Entropy.exists_canonicalHalfSampleTrimIterate_joint_probability} in
\doclink{QICLean/Analysis/StoppedHalfSamplingProbability}.}
This establishes a mathematical probability statement for the corrected
repeated-resampling operation under its stated family and reference
hypotheses. It does not identify that operation with the printed
sample-once procedure or establish running time.

\bibliographystyle{alpha}
\bibliography{references}
\end{document}
